\section*{3. 1.3 的证明}

\subsubsection*{3.1.}
在定理的假设下, 只须证明 1.5(b) 不成立. $h^{0,2n}\ne0$ 的情形将由 Katz 在 XXI 4.1 中处理（条件 $h^{0,2n}\ne0$ 等价于
\[
  \sum_i a_i-(2n+r+1)\geq0,
\]
此数就是满足 $\Omega_V^{2n}\simeq\mathcal O(m)$ 的整数 $m$）.
\footnote{原印本在此和式中写作 $d_i$, 虽然多重次数记为 $(a_i)$.}
本讲的结果在讨论班余下部分不再使用, 因而不存在循环论证.

\subsubsection*{3.2.}
当 $k=\mathbb C$ 时, 证明 1.3 只须注意, 若 1.5(b) 成立, 则 $V$ 的上同调纯为 $(n,n)$ 型, 因为代数上同调为 $(n,n)$ 型. 这只在 1.3 所排除的例外情形中发生（XI 2.9）. 因此由 Lefschetz 原理, 1.3 在特征零时成立.

\subsubsection*{命题 3.3.}
采用 1.4 的假设, 并假设 (b) 成立. 设 $\operatorname{Ev}(V,\mathbb Q)$ 为 $\operatorname{Ev}(V,\mathbb Q_\ell(n))$ 中由代数上同调类生成的 $\mathbb Q$-向量子空间.

\begin{enumerate}[label=(\alph*)]
\item 有同构
  \[
  \operatorname{Ev}(V,\mathbb Q)\otimes\mathbb Q_\ell
  \xrightarrow{\ \sim\ }
  \operatorname{Ev}(V,\mathbb Q_\ell(n)).
  \]
  在 $\operatorname{Ev}(V,\mathbb Q)$ 上, 相交型是有理且定号的, 符号为 $(-1)^n$.
\item 消灭圈及其共轭属于 $\operatorname{Ev}(V,\mathbb Q)$; 它们在那里组成 $A$, $D$ 或 $E$ 型的不可约根系.
\item 单值群（即 $\pi_1(\ell-S,\bar s)$ 在 $\operatorname{GL}(\operatorname{Ev}(V,\mathbb Q))$ 中的像）为有限群; 它就是 (b) 中根系的 Weyl 群.\footnote{原印本反复使用“groupe de Weil”; 根系所对应的群是其 Weyl 群.}
\end{enumerate}

除相交型的定号性外, (a) 由相交型在 $\operatorname{Ev}(V,\mathbb Q_\ell(n))$ 上非退化（由 (LV)）, 且在 $\operatorname{Ev}(V,\mathbb Q)$ 上有理推出.

条件 1.4(b) 等价地表示, $V$ 上存在模数值等价后的代数圈, 它们在单值作用下不保持不变; 因此此条件与 $\ell$ 无关. 设 $x$ 为这样的代数圈, $\delta$ 为由 $s_0\in S$ 以及从 $\bar s$（它定义 $V$）到 $s_0$ 的路径 $c$ 所对应的消灭圈. 对适当的 $c$, 有 $(x,\delta)\ne0$. Galois 变换
\[
  \sigma(x)=x-(-1)^n(x,\delta)\delta
\]
仍为代数的; 因此 $(x,\delta)\delta$ 是代数的, 它是一个与 $\ell$ 无关的圈 $y$ 的像. 所以有理数
\[
  \frac{(-1)^n}{2}(y,y)=(x,\delta)^2
\]
与 $\ell$ 无关. 它在每个 $\mathbb Q_\ell$ 中都是平方, 因而在 $\mathbb Q$ 中也是平方; 因此, 不计符号, $(x,\delta)$ 属于 $\mathbb Q$ 且与 $\ell$ 无关. 这不仅说明 $\delta\in\operatorname{Ev}(V,\mathbb Q)$, 还说明模数值等价的代数圈
\[
  \pm\delta=\frac{\sigma(x)-x}{(x,\delta)}
\]
与 $\ell$ 无关.

簇 $V$ 上的代数圈, 总是在 $V$ 的任一定义域的某个有限扩张上有定义. 因此 $\pi_1(\ell-S,s)$ 通过一个有限群作用于 $\operatorname{Ev}(V,\mathbb Q)$. 此群由对应于消灭圈 $\delta$ 及其共轭的反射生成（XVIII 6.6.1 及第 2 节）. 因此各圈 $\delta$ 及其共轭组成一个根系, 其 Weyl 群就是单值群. 此外, 这个根系的所有根在 Weyl 群下共轭; 它为 $A$, $D$ 或 $E$ 型.

在相差一个标量的意义下, 相交型是 Weyl 群下唯一的不变型. 因而它是定号的. 其符号由计算
\[
  (\delta,\delta)=2(-1)^n.
\]
确定.

\subsubsection*{命题 3.4.}
若 $k=\mathbb C$, 则在 1.4 的假设下, 下列条件等价:
\begin{enumerate}[label=(\roman*)]
\item 单值群有限;
\item 消灭上同调 $\operatorname{Ev}(V)$ 纯为 $(n,n)$ 型.
\end{enumerate}

设 $\widetilde X\to\ell$ 为纤维等于线束 $\ell$ 中超平面截面的族, $V_s$ 为 $s\in\ell(\mathbb C)$ 处的纤维. 各个 $\operatorname{Ev}(V_s)$ 在 $\ell$ 上组成 Hodge 结构的代数变分. 对 $\ell$ 的一个有限覆盖应用固定部分定理 [2, 4.1.2]（参见 [2, 4.1.3.3]）, 由 (i) 可知 $\operatorname{Ev}(V_s)$ 上的 Hodge 结构局部常值. 因此单值群通过 Hodge 结构的自同构作用. 对一个 Picard--Lefschetz 变换 $T$, $\operatorname{Im}(T-1)$ 是秩一的 Hodge 子结构, 因而为 $(n,n)$ 型, 且由一个消灭圈 $\delta$ 生成. 因为消灭圈为 $(n,n)$ 型, 所以整个 $\operatorname{Ev}(V_s)$ 都为此型.

反之, 若 $\operatorname{Ev}(V)$ 纯为 $(n,n)$ 型, 则 $\operatorname{Ev}(V)$ 上的相交型定号, 单值群有限（参见 [2, 4.1.3.3]）.

\subsubsection*{注 3.5.}
若 3.4 中的等价条件成立,\EditorialFootnote{法文印本及英文工作版此处均写作 3.3; 等价条件实际列于命题 3.4.} 则 3.3 的论证说明, 消灭圈及其共轭组成 $A$, $D$ 或 $E$ 型根系, 其 Weyl 群就是单值群.

\subsubsection*{3.6.}
证明当 $k$ 的特征为不同于 $2$ 的有限特征时 1.3 成立. 设 $S=\operatorname{Spec}(\mathcal V)$ 为剩余域 $k$ 的混合特征离散赋值环的谱, 并设 $X\subset\mathbb P_S^{2n+d}$ 为 $S$ 上相对维数 $2n+1$、多重次数 $(a_1,\ldots,a_{d-1})$ 的完全交. 设 $\ell$ 为次数 $a_d$ 的超曲面在 $X$ 上截出的截面的 Lefschetz 线束（在特殊纤维中, $\ell_s$ 是 $X_s$ 的 Lefschetz 线束）. 设
\[
  f:\widetilde X\longrightarrow\ell\simeq\mathbb P^1_S
\]
为纤维等于此线束中各截面的族. 例外轨迹 $T\subset\ell$ 在 $S$ 上平展（XVII 6; 此处使用 $p\ne2$）. 层 $R^{2n}f_*\mathbb Q_\ell$ 在 $\ell-T$ 上局部常值; 因而它温和分歧, 且其在特征 $p$ 时的单值作用与在特征零时相同. 若 1.5(b) 成立, 则由 3.3(c) 与 3.4 可知, 维数 $2n$、多重次数 $\underline a$ 的复完全交的上同调群 $H^{2n}(V)$ 纯为 $(n,n)$ 型. 这只在 1.3 所排除的情形中发生（XI 2.9）.
